\section{Формулы. Реализация функций формулами, интерпретация формул. Таблицы истинности. Равносильные формулы.}

\subsection*{Формулы и их реализация}

\paragraph{Определение формулы}
\textbf{Формула} над множеством функций $F$ (называемым базисом) — это цепочка символов, построенная по следующим правилам:
\begin{enumerate}
    \item Любая переменная $x_i$ является формулой.
    \item Если $\mathcal{J}_1, \dots, \mathcal{J}_n$ — формулы, а $f \in F$ — функция арности $n$, то выражение $f(\mathcal{J}_1, \dots, \mathcal{J}_n)$ является формулой.
\end{enumerate}
В записи часто используют инфиксную форму ($a \lor b$ вместо $\lor(a, b)$) и приоритет операций ($\neg, \&, \lor, \to$), чтобы опускать лишние скобки.

\paragraph{Интерпретация (Алгоритм Eval)}
Каждой формуле $\mathcal{J}$ соответствует булева функция. Процесс вычисления значения формулы на наборе значений переменных называется \textbf{интерпретацией}.
\textbf{Рекурсивный алгоритм Eval:}
\begin{itemize}
    \item Если $\mathcal{J}$ — переменная $x_i$, вернуть её заданное значение.
    \item Если $\mathcal{J} = f(\mathcal{J}_1, \dots, \mathcal{J}_n)$, то сначала вычислить значения всех подформул $\mathcal{J}_i$, а затем применить к ним операцию $f$.
\end{itemize}
Если формула $\mathcal{J}$ при интерпретации на всех наборах выдает те же значения, что и функция $f$, говорят, что \textbf{формула реализует функцию} ($\text{func } \mathcal{J} = f$).

\subsection*{Таблицы истинности и равносильность}

\paragraph{Построение таблицы}
Для определения функции, которую реализует формула, строится таблица истинности. Для этого формула разбивается на подформулы (операции), и для каждой вычисляется промежуточный столбец.
\textit{Пример:} Для формулы $\mathcal{J} = (x_1 \land x_2) \lor (\neg x_1 \land x_2)$ итоговый столбец совпадет со столбцом функции $x_2$.

\paragraph{Равносильные формулы}
Две формулы $\mathcal{J}_1$ и $\mathcal{J}_2$ называются \textbf{равносильными} ($\mathcal{J}_1 \equiv \mathcal{J}_2$), если они реализуют одну и ту же функцию.
Отношение равносильности является отношением эквивалентности.

\paragraph{Основные равносильности (Законы булевой алгебры)}
Новиков выделяет 10 базовых групп (проверяются таблицами истинности):
\begin{enumerate}
    \item \textbf{Идемпотентность:} $a \lor a = a, a \land a = a$.
    \item \textbf{Коммутативность:} $a \lor b = b \lor a, a \land b = b \land a$.
    \item \textbf{Ассоциативность:} $a \lor (b \lor c) = (a \lor b) \lor c$.
    \item \textbf{Поглощение:} $(a \land b) \lor a = a, (a \lor b) \land a = a$.
    \item \textbf{Дистрибутивность:} $a \land (b \lor c) = (a \land b) \lor (a \land c)$.
    \item \textbf{Константы:} $a \lor 1 = 1, a \land 0 = 0$.
    \item \textbf{Нейтральные элементы:} $a \lor 0 = a, a \land 1 = a$.
    \item \textbf{Двойное отрицание:} $\neg \neg a = a$.
    \item \textbf{Законы де Моргана:} $\neg(a \land b) = \neg a \lor \neg b, \neg(a \lor b) = \neg a \land \neg b$.
    \item \textbf{Дополнение:} $a \lor \neg a = 1, a \land \neg a = 0$.
\end{enumerate}

\subsection*{Подстановка и замена}

\paragraph{Правило подстановки}
Если в равносильной формуле заменить все вхождения некоторой переменной $x$ на одну и ту же формулу $\mathcal{A}$, то полученные формулы останутся равносильными. Это позволяет выводить сложные равносильности из простых.

\paragraph{Правило замены}
Если в формуле $\mathcal{A}$ заменить некоторую подформулу $\mathcal{B}$ на равносильную ей подформулу $\mathcal{C}$ ($\mathcal{B} \equiv \mathcal{C}$), то полученная формула $\mathcal{A}'$ будет равносильна исходной ($\mathcal{A} \equiv \mathcal{A}'$).

\paragraph{Примеры}
\begin{enumerate}
    \item Формула $((x_1 \land x_2) \to x_1)$ реализует \textbf{константу 1} (тавтология).
    \item Формула $((x_1 \land x_2) + x_1) + x_2$ реализует \textbf{дизъюнкцию} $x_1 \lor x_2$.
\end{enumerate}